%% =============================================================================
%% COMPARISON AND VALIDATION
%% =============================================================================

\section{Comparison and validation} \label{sec:comparison}

A statistical software package must demonstrate both correctness and competitive
performance. This section compares \pkg{prompt2analytics} against established
\proglang{R} packages on both dimensions, showing that the \proglang{Rust}
implementation achieves speedups ranging from 2$\times$ to over
100$\times$, depending on the method, while matching reference implementations
to high numerical precision.

\subsection{Methodology}

Comparisons use synthetic datasets generated with identical data-generating
processes in both languages (seed $= 42$) and classic datasets with published
results (Longley for OLS, Grunfeld for panel data). Reference implementations
include \fct{lm}/\pkg{sandwich} for OLS, \pkg{plm}/\pkg{lfe} for panel data,
\fct{glm} for discrete choice, and \pkg{forecast} for time series. Each
benchmark executes 100 iterations after warmup; all tests run on identical
hardware (Intel i7-1260P, 64GB RAM). Appendix~\ref{app:benchmark-methodology}
provides complete details on data generating processes, execution protocol,
and software versions.

\subsection{Performance comparison}

Table~\ref{tab:performance-summary} summarizes performance results across
representative methods and sample sizes. The speedup factors vary considerably
by method category and sample size, reflecting differences in algorithmic
complexity and the relative weight of interpreter overhead versus numerical
computation.

\begin{table}[htbp]
\centering
\caption{Performance comparison: median execution time ($\mu$s) and speedup
factors for representative methods}
\label{tab:performance-summary}
\renewcommand{\arraystretch}{1.2}
\small
\begin{tabular}{llrrrr}
\toprule
Category & Method & $n$ & \proglang{R} & \pkg{p2a} & Speedup \\
\midrule
\multirow{3}{*}{Regression} & OLS & 100 & 812 & 42 & 19$\times$ \\
 & OLS & 10,000 & 2,181 & 924 & 2.4$\times$ \\
 & OLS + HC1 & 10,000 & 25,224 & 1,571 & 16$\times$ \\
\midrule
\multirow{2}{*}{Panel} & Fixed effects & 100 & 4,901 & 27 & 184$\times$ \\
 & HDFE (2-way) & 5,000 & 26,657 & 1,159 & 23$\times$ \\
\midrule
\multirow{2}{*}{Discrete} & Logit & 1,000 & 2,176 & 401 & 5.4$\times$ \\
 & Probit & 1,000 & 2,760 & 2,347 & 1.2$\times$ \\
\midrule
\multirow{2}{*}{Time series} & ARIMA(1,1,1) & 500 & 4,783 & 521 & 9.2$\times$ \\
 & MSTL & 500 & 1,853 & 219 & 8.5$\times$ \\
\midrule
\multirow{2}{*}{ML} & K-means ($k=3$) & 1,000 & 1,274 & 587 & 2.2$\times$ \\
 & PCA & 1,000 & 246 & 67 & 3.7$\times$ \\
\bottomrule
\end{tabular}
\end{table}

Several patterns emerge from these results. First, speedup factors are largest
for methods with significant per-call overhead in \proglang{R}. Panel data
methods show the largest improvements: fixed effects estimation at
$n = 100$ is 184$\times$ faster in \pkg{prompt2analytics}. This extreme
speedup is misleading for practical interpretation: it primarily reflects
\proglang{R}'s interpreter overhead (formula parsing, factor handling, data frame
manipulation) rather than differences in numerical computation. At small sample
sizes, overhead dominates; the same matrix algebra that takes microseconds is
wrapped in milliseconds of bookkeeping. As sample size increases and matrix
operations dominate execution time, speedup factors converge to more realistic
levels: 5--6$\times$ at $n = 10{,}000$. Users should focus on the
$n = 10{,}000$ results as indicative of typical production performance. The
$n = 100$ results demonstrate overhead elimination but not sustained computational
advantage.

Second, methods involving expensive mathematical functions show smaller
speedups. Probit estimation requires repeated evaluation of the standard
normal CDF and PDF; both \proglang{R} (via compiled C code) and
\proglang{Rust} incur similar costs for these operations, yielding a
modest 1.2$\times$ speedup at $n = 1{,}000$. Logit, which uses the
computationally simpler logistic function, achieves 5.4$\times$ speedup
at the same sample size.

Third, robust standard error computation shows consistent speedups (13--16$\times$)
across sample sizes. The HC1 estimator requires computing the ``meat'' matrix
$\sum_i \hat{u}_i^2 x_i x_i'$, which benefits from \proglang{Rust}'s efficient
memory access patterns and SIMD vectorization in the \pkg{faer} linear algebra
library.

Figure~\ref{fig:benchmark-boxplots} displays the full timing distributions for
selected methods at $n = 10{,}000$. Beyond the differences in central tendency,
\pkg{prompt2analytics} shows substantially lower variance: interquartile ranges
are typically 10--100$\times$ smaller than in \proglang{R}, reflecting the
absence of garbage collection pauses and interpreter overhead. This predictability
is valuable for production deployments where consistent response times matter.

\begin{figure}[htbp]
\centering
\includegraphics[width=\textwidth]{figures/benchmark_boxplots}
\caption{Execution time distributions for \proglang{R} (blue) and
\pkg{prompt2analytics} (orange) at $n = 10{,}000$. Y-axis uses logarithmic
scale. Boxes show interquartile range; whiskers extend to 1.5$\times$ IQR.}
\label{fig:benchmark-boxplots}
\end{figure}

Memory usage differences are also notable. For OLS at $n = 10{,}000$,
\proglang{R} allocates approximately 3.6~MB (measured via \fct{bench::mark})
while \pkg{prompt2analytics} reports minimal allocation, under 1~KB beyond
the pre-allocated data structures, because operations proceed in-place where
possible and temporary buffers are reused across iterations. For robust
standard errors at the same sample size, \proglang{R} allocates over 10~MB
due to intermediate matrix copies; \pkg{prompt2analytics} avoids these
through careful ownership management. \proglang{Rust} allocations are
measured via custom instrumentation of the global allocator during
benchmark runs.

\subsection{Numerical validation}

We validate all methods against reference implementations using sample-size-dependent
tolerances (looser for small samples where ill-conditioning amplifies differences,
tighter for large samples). The Longley dataset, with its extreme multicollinearity
(condition number $> 10^9$), provides a stringent test. Table~\ref{tab:validation-ols}
shows coefficients match \proglang{R}'s \fct{lm} to at least 8 significant digits.

\begin{table}[htbp]
\centering
\caption{OLS validation on the Longley dataset ($n = 16$, $k = 6$)}
\label{tab:validation-ols}
\renewcommand{\arraystretch}{1.15}
\begin{tabular}{lrrr}
\toprule
Coefficient & \proglang{R} \fct{lm} & \pkg{prompt2analytics} & Difference \\
\midrule
(Intercept) & $-3482.259$ & $-3482.259$ & $< 10^{-8}$ \\
GNP.deflator & $0.01506$ & $0.01506$ & $< 10^{-8}$ \\
GNP & $-0.03582$ & $-0.03582$ & $< 10^{-8}$ \\
Unemployed & $-0.02020$ & $-0.02020$ & $< 10^{-8}$ \\
Armed.Forces & $-0.01033$ & $-0.01033$ & $< 10^{-8}$ \\
Population & $-0.05110$ & $-0.05110$ & $< 10^{-8}$ \\
Year & $1.82915$ & $1.82915$ & $< 10^{-8}$ \\
\midrule
$R^2$ & $0.9955$ & $0.9955$ & $< 10^{-4}$ \\
\bottomrule
\end{tabular}
\end{table}

Heteroskedasticity-consistent standard errors (HC0--HC3) are validated against
the \pkg{sandwich} package. For HC0 and HC1, which differ only by a degrees-of-freedom
correction, results match to $< 10^{-8}$. For HC2 and HC3, which involve leverage
values $h_{ii} = x_i'(\mathbf{X}'\mathbf{X})^{-1}x_i$, results match to $< 10^{-7}$.
The slightly larger tolerance reflects algorithmic differences in leverage
computation: \pkg{sandwich} extracts leverage from the QR decomposition used
internally by \fct{lm}, while \pkg{prompt2analytics} computes the hat matrix
diagonal directly via matrix multiplication. Both approaches are numerically
valid; the discrepancy arises from different accumulation of floating-point
rounding errors rather than any methodological difference. For applied use,
differences at the $10^{-7}$ level are negligible.

Panel data methods are validated using the Grunfeld dataset (10 firms, 20 years,
$n = 200$). Fixed effects coefficients match \pkg{plm} to $< 10^{-6}$, and the
within $R^2$ matches to four decimal places. The Hausman test statistic for
choosing between fixed and random effects matches to $< 10^{-4}$. Two-way
high-dimensional fixed effects, implemented using the MAP
algorithm~\citep{Gaure:2013}, matches \pkg{lfe}'s \fct{felm} to $< 10^{-5}$.

Beyond these flagship methods, the validation suite covers over 200 methods
with more than 1,500 individual test cases (334 specifically targeting validation
against R reference implementations). Statistical tests ($t$-tests, ANOVA,
chi-squared, Fisher's exact, Wilcoxon, Shapiro-Wilk, Kolmogorov-Smirnov) are
validated against \proglang{R}'s \pkg{stats} package. Discrete choice models
(logit, probit) match \fct{glm} coefficients and standard errors. Time series
methods (VAR, ARIMA) match the \pkg{vars} and \pkg{forecast} packages.
Clustering algorithms match \fct{kmeans} cluster assignments and within-cluster
sum of squares. All validation tests are integrated into continuous integration
and can be run via \code{cargo test -p p2a-core}.

\subsection{Sources of performance differences}

The performance advantages of \pkg{prompt2analytics} stem from several
architectural choices in \proglang{Rust} and the library design.

First, \proglang{Rust}'s ownership system enables stack allocation and in-place
mutation without garbage collection. Matrix operations can reuse memory rather
than allocating new buffers, and there are no GC pauses during computation.
This is particularly impactful for iterative algorithms (HDFE, MLE) where
\proglang{R} would allocate temporary objects on each iteration.

Second, the \pkg{faer} linear algebra library provides highly optimized
implementations that leverage SIMD instructions on modern processors. While
\proglang{R} with OpenBLAS also uses SIMD for matrix multiplication, the
overhead of crossing the R-C boundary for each operation accumulates.

Third, the column-based API eliminates formula parsing. \proglang{R}'s model
formula interface (\code{y \~{} x1 + x2}) involves parsing, factor expansion,
and model matrix construction---operations that can dominate execution time
for small datasets. \pkg{prompt2analytics} receives pre-specified column names
directly from the MCP tool call.

Fourth, the single-language design avoids the overhead of data serialization
between components. In typical \proglang{R} workflows, data moves between
\proglang{R}, C, and potentially Fortran; each transition involves memory
copies and type conversions. \pkg{prompt2analytics} maintains data in native
\proglang{Rust} structures throughout.

These advantages diminish as sample size increases and matrix algebra dominates.
For very large datasets ($n > 100{,}000$), speedup factors converge toward
2--3$\times$ as both implementations spend most time in optimized BLAS routines.
The largest advantages appear for small-to-medium datasets ($n = 100$ to $10{,}000$)
and for methods with significant per-call overhead.

\subsection{Feature comparison}

Performance is one dimension of software quality; feature completeness and
ecosystem integration are equally important for applied work. Table~\ref{tab:feature-comparison}
compares \pkg{prompt2analytics} against established \proglang{R} packages across
dimensions beyond execution speed.

\begin{table}[htbp]
\centering
\caption{Feature comparison with established \proglang{R} packages}
\label{tab:feature-comparison}
\small
\renewcommand{\arraystretch}{1.15}
\begin{tabular}{@{}lcccc@{}}
\toprule
Feature & \pkg{p2a} & \pkg{fixest} & \pkg{lfe} & \pkg{plm} \\
\midrule
\multicolumn{5}{l}{\textit{Panel data methods}} \\
Within (FE) estimator        & \cmark & \cmark & \cmark & \cmark \\
Random effects (GLS)         & \cmark & ---    & ---    & \cmark \\
High-dimensional FE          & \cmark & \cmark & \cmark & --- \\
Multi-way clustering         & 2-way  & $n$-way & 2-way & 2-way \\
Hausman test                 & \cmark & ---    & ---    & \cmark \\
\addlinespace
\multicolumn{5}{l}{\textit{Standard errors}} \\
HC0--HC3 robust              & \cmark & \cmark & \cmark & \cmark \\
Driscoll-Kraay               & \cmark & \cmark & ---    & \cmark \\
Newey-West HAC               & \cmark & \cmark & \cmark & \cmark \\
Conley spatial               & ---    & \cmark & \cmark & --- \\
\addlinespace
\multicolumn{5}{l}{\textit{Causal inference}} \\
Difference-in-differences    & \cmark & \cmark & --- & --- \\
Callaway-Sant'Anna (staggered) & \cmark & \cmark & --- & --- \\
IV/2SLS                      & \cmark & \cmark & \cmark & \cmark \\
Weak instrument tests        & F-stat & \cmark & \cmark & --- \\
\addlinespace
\multicolumn{5}{l}{\textit{Diagnostics \& output}} \\
Coefficient plots            & ---    & \cmark & --- & --- \\
Regression tables (export)   & JSON   & \LaTeX/HTML & \LaTeX & \LaTeX \\
Marginal effects             & AME    & \cmark & --- & --- \\
Formula interface            & ---    & \cmark & \cmark & \cmark \\
\addlinespace
\multicolumn{5}{l}{\textit{Integration}} \\
LLM/MCP interface            & \cmark & ---    & ---    & --- \\
Local execution (no R)       & \cmark & ---    & ---    & --- \\
tidyverse compatible         & ---    & \cmark & ---    & \cmark \\
\bottomrule
\end{tabular}
\end{table}

For core panel methods, \pkg{prompt2analytics} provides comparable coverage
to specialized packages; the primary remaining gap is Conley spatial standard
errors. The distinguishing features are the natural language interface and
fully local execution without \proglang{R} dependencies. Users requiring
maximum performance for high-dimensional fixed effects should note that
\pkg{fixest} likely outperforms \pkg{prompt2analytics} due to its highly
optimized C++ implementation specifically tuned for this problem class.

Performance advantages become practically relevant for simulation
studies, bootstrap procedures, and deployment contexts where \proglang{R}
installation is impractical. For typical one-off analyses, speed differences
are imperceptible.

\subsection{Limitations}

Several caveats apply to these comparisons. First, we compare against specific
\proglang{R} packages chosen for their widespread use in applied econometrics;
the comparison necessarily reflects a snapshot of rapidly evolving software.
Second, benchmarks use synthetic data with simple structure; real datasets
with missing values, factors, or complex hierarchies may show different
relative performance. Third, \proglang{R}'s flexibility (formula interface,
extensive diagnostics, integration with the tidyverse ecosystem) provides
value not captured in raw timing comparisons or feature checklists.

\subsubsection{Comparison with fixest}

\pkg{fixest}~\citep{Berge:2018} was excluded from the performance benchmarks
in Table~\ref{tab:performance-summary} because its highly optimized C++
implementation makes direct speed comparison uninformative for most users.
\pkg{fixest} is specifically engineered for high-dimensional fixed effects
and will outperform \pkg{prompt2analytics} on HDFE problems, often by a
substantial margin.

Users who prioritize raw speed for high-dimensional fixed effects estimation
should use \pkg{fixest} directly---this is the right tool for that job.
The differentiating value of \pkg{prompt2analytics} is not speed leadership
on any single method class, but rather: (1)~the natural language interface
that enables users unfamiliar with \proglang{R} syntax to access validated
methods, (2)~local execution without \proglang{R} dependencies, and
(3)~integration across a broad range of econometric methods within a
unified framework.

For users already proficient in \proglang{R} who work primarily with
high-dimensional fixed effects, \pkg{fixest} remains the recommended choice.
\pkg{prompt2analytics} complements rather than replaces specialized tools.

The performance advantages documented here become practically relevant
for interactive use, where sub-second response times matter for
conversational workflows. Section~\ref{sec:performance} presents
detailed benchmarks across approximately 75 method configurations.

